\section{Why the logarithmic-order breakthrough is still missing}
\label{sec:limits}

The new construction changes the finite contraction menu and its interaction
with the global recursion. It reduces the coefficient multiplying
\[
 K(M,u)=\frac{\pi\sqrt2}{\ln2}\frac{M}{\sqrt u}.
\]
It does not change the dependence on $u$. The admissible root taper remains
$u=\Theta(\log s)$, so the normalized order remains
$s/\sqrt{\log s}$. The occurrence bridge then retains that denominator.
The compressed reduction still incurs an $O(n\log\log n)$ incidence
budget at unrestricted density~\cite{upstream,degree}.

\subsection{The useful new idea and its limitation}

Partial duplication deliberately trades a slightly larger local time
coefficient for a better global choice of taper and child shrinkage.
The fifteen-group menu has local $\beta_{15}\approx0.57154$, whereas
the full eighteen-group menu has $\beta_{18}\approx0.56529$.
The smaller local number alone does not determine the best root bound.
For fifteen groups the root taper is $0.1134\lambda$, compared with
$0.107\lambda$ for eighteen groups; the recursive entries remain cheap
enough after shrinking by $2/15$. The resulting root coefficient is smaller
even though its local norm certificate is larger.

Uniformly cloning still more points is not a route to a smaller order
within this certificate. If every affine point has multiplicity $t$, then
$r=9t$. Final-block support self norms and cross products both scale as
$t^2$, while the entropy normalization divides by $r^2=81t^2$.
Their normalized contribution therefore stays fixed. More explicitly,
if $y=(1,47/56+\varepsilon,15/28)$, the final-operation coupled bound
already contributes
\[
 \frac1{18}\sum_i y_i^2+\frac1{27}\left(\sum_i y_i\right)^2
\]
to $\beta^2$, independently of $t$. Internal operations can only increase
the maximum used by this certificate. Meanwhile the sufficient packing
capacity decreases as $1/(1+2\log_2r)$. Increasing a fixed cloning
multiplicity therefore does not make the coefficient vanish. Letting $r$
grow with input size would additionally require new uniform allocation,
padding and enumeration bounds; the fixed-$r$ analysis cannot be reused
without those charges.

\subsection{Routes checked against the latest sources}

\paragraph{Decision-only precision.}
Parity, isolation or reducing the requested output precision does not
remove the current reconstruction engine's internal
$\Theta(MD/u)$ precision loss. A faster decision procedure would have to
avoid that recovery step with a new exact summation or reconstruction
identity. This obstacle already appears in the asymptotic-analysis and
precision-refinement packets~\cite{barriers,precision}.

\paragraph{A thinner graded space with the same crossing decomposition.}
The latest weighted-pair packet proves a sharper obstruction than merely
counting oscillator states~\cite{pairwise}. Under its full-crossing-rank
and polynomial-piece assumptions, a crossing on an $N$-symbol alphabet
factors through energy windows $E\le4h$ only if
\[
 N^2\le\poly(h)\max_{E\le4h}\dim W_E.
\]
If $\log_2N\sim ch^\alpha$ and
$\log_2\dim W_E\sim cE^\alpha$ with the required monotonicity, then
\[
 2ch^\alpha\le c(4h)^\alpha+o(h^\alpha),
 \quad\hbox{hence}\quad \alpha\ge\tfrac12.
\]
Thus simply proposing $\alpha<1/2$ while retaining all those interfaces
is inconsistent. The obstruction is scoped: a different energy window,
a different reconstruction identity, a restricted image, or structured
summation could fall outside it. It does not prove a lower bound for
all algorithms deciding 3-SAT.

\paragraph{Faster matrices, more groups, or more trees.}
These can improve constants, but complete explicit boundary tables still
pay their array sizes. A larger menu cannot remove the cost of materializing
the tables it actually uses. Existing source analyses already isolate
this obstruction for the full-table model~\cite{tree,pairwise}.
Likewise, increasing the group count in the dense reduction improves its
degree premise without by itself removing the incidence-indexing budget.
The density-sensitive and high-moment proposals already exploit that
possibility~\cite{degree,project}.

\subsection{A concrete next target beyond coefficient optimization}

The most useful next theorem would concern the \emph{actual family} of
deformed group factors generated from 3-SAT equality and clause tensors.
The present entropy relaxation treats their boundary data as arbitrary
arrays of the permitted sizes. A strict improvement in order needs a
representation or summation rule that uses additional algebraic structure.
One precise research target is a family of exact boundary representations
closed under the required pair contractions and recovery operations, with
construction and evaluation cost
\[
 2^{o(M/\sqrt u)}\,\poly(s)
\]
at the critical root scale. A certificate must include the representation
dimension, coefficient precision, update costs and recursive entry costs.
A compact description whose evaluation still enumerates all boundary
indices does not meet this target.

This suggests a falsifiable program: characterize the image spaces of
small genuine 3-SAT group factors after deformation; test whether a
proposed exact representation survives one merge and recovery step;
then prove the uniform closure and cost bounds. Full matrix rank rules out
some factor-through-small-space proposals immediately, but does not itself
rule out fast structured summation. No such closed representation theorem
has been constructed in this report. The existing conditional
$2^{O(n\log\log n/\log n)}$ target still depends on its unproved
structured-mode acceleration hypothesis~\cite{precision}.

\section{Verification, reproducibility, and remaining obligations}
\label{sec:verification}

The archive contains complete LaTeX sources, the compiled report, exact
standard-library Python checkers, deterministic JSON results, source and
commit ledgers, independent audit notes, and a checksum manifest.
There are two independently reconstructed certificates for the main
fifteen-group geometry, with different point-class implementations.

\begin{enumerate}
\item \path{audit_round5_new15.py} constructs all four root plans and
their $117$ supports, checks $82{,}134$ cross-plan time pairs, and proves
the strict schedule inequalities with rational interval arithmetic.
It certifies all four fully paid branches and the time and space radical
identities. The ten-, seven-, and four-group local bounds are explicitly
imported source inputs.
\item \path{check_partial_affine_joint.py} independently constructs the
geometry and verifies $82{,}134$ time pairs, $4{,}374$ boundary pairs and
$37{,}908$ mixed pairs. It implements the rounded integer joint selector
and independently verifies the complete time and space schedule.
\item \path{round5_affine_certificate.py} and
\path{check_rectangular_joint.py} verify the eighteen-group alternative.
The latter checks $124{,}416$ time pairs, $6{,}534$ boundary pairs and
$57{,}024$ mixed pairs. The former also checks the nine-group rectangular
menu and a matching witness for the eighteen-group \emph{local} Euclidean
certificate; it gives no 3-SAT lower bound.
\item \path{audit_round5_baseline.py} independently reconstructs the
prior cyclic thirteen-plan certificate with the compatible rectangular
backend and checks the stronger $13\to10\to7\to6\to4$ comparison.
The exact source coefficients imported at the remaining arities are
identified in the audit note and dependency ledger.
\item \path{check_tensor_regrouping.py} compares literal directed-index
summation with every affine plan in $36$ fixtures, checking $144$ complete
plan contractions and $24{,}192$ literal index assignments over
$(\mathbb Z/256\mathbb Z)[V]/(V^2)$. It includes asymmetric directed
dimensions, zero factors, nonzero cases and two nilpotent factors, and
checks the actual integer selectors on the fifteen- and eighteen-group
fixtures. These are algebra tests, not fast-matrix benchmarks.
\item \path{verify_occurrence_bridge.py} compares Boolean assignment
counts with incidence-wire sums on $483$ small formula cases and
$1{,}293$ threshold/branch comparisons. It tests unused variables,
tautologies, repeated literals, empty clauses and satisfied clauses whose
remaining incidences must be retained.
\end{enumerate}

The coefficient enclosures use Machin's formula with alternating-series
remainders for $\pi$, a positive atanh series and geometric tail for
$\ln2$, and integer-square-root radical bounds. Floating-point numbers
are used only for explicitly exploratory displays. All delivered main
checks use explicit exceptions and remain active under Python's
optimization flag. Ordinary and optimized replay both pass.

The finite tests establish finite identities and certificates. The proof
of uniform asymptotic complexity is the schedule argument, together with
the imported analytic interface. No full fast SAT solver, asymptotic
matrix implementation, or new Lean/Rocq build is supplied or claimed.
The finite nearby-menu exploration is recorded separately as exploratory
evidence; it is not an optimality proof and is not needed by either theorem.

\noindent\begin{minipage}{\linewidth}
From the extracted packet directory, run:
\begin{lstlisting}
python3 verify_manifest.py
python3 run_checks.py
python3 run_checks.py --optimized
python3 build_report.py
\end{lstlisting}
\end{minipage}

The checkers use only the Python standard library. Building the report
requires pdfLaTeX and the standard packages listed in the preamble; shell
escape is disabled. Replay rewrites the deterministic result files.
The manifest should be checked before rebuilding the PDF, whose metadata
may depend on the TeX environment. The report was compiled and visually
inspected before delivery.
